% Transcription — fonds Alexandre Grothendieck, shelfmark 36, batch 1 (pages 1-20).
% Established from the facsimile archives/batches/36/batch-01.pdf (a copy of
% 36.pdf fetched through the project's relay, grothendieck-relay.onrender.com;
% PDF page k = archivists' page k, no cover sheet), per the skill
% transcribe-grothendieck. The folder is 84 pages in five batches; this is the
% first (pages 21-84 are batches 2-5, done separately).
%
% Pass: Opus 5.5 (claude-opus-5-5), 2026-09-26 — first pass, unchecked against the pages by a human.
%
% Pages skipped: 1 (a blank cream sheet carrying only show-through of the
% squared leaf behind it) and 8 (blank). Pages noted only (another hand, not
% re-set): 2-7. Pages transcribed: 9-20.
%
% Pages 2-7: six leaves of squared paper in blue ink, numbered 1-6 at the top
% right, headed « Remarques. SGA 7 IX »: an unidentified reader's remarks on a
% typed state of Exposé IX of SGA 7 (« Modèles de Néron et monodromie », his
% exposé), keyed to its sections (circled 1-11, 8 left empty) and to the
% typescript's pages (0-103) and lines. The writer addresses the exposé's
% author as « tu » (« As tu vérifié si l'accouplement d'A.M. coïncide avec le
% tien ? »), so the ink is not taken as his. By the project rule (as for
% Tate's notes, folder 6 batch 2, p. 38) another person's text is not re-set:
% each leaf gets \page{N} and ONE French \note identifying it (leaf n of 6,
% sections and typed pages covered, a one-line summary of the mathematics),
% then giving in full, with their place on the page, the replies attributed
% to the exposé's author — pencil: « guillemets suffisants ? », « je le
% dis ! », « c'est dit dans la 5.13 », « c'est pourtant trivial », « non, mon
% argument marche », « que faire ? » (twice), « si ? », question marks and
% circles; ink: « CANULAR ?! » — said to be presumably his hand, not
% established. (Revised at the coordinator's request; a first version re-set
% the remarks in full.)
%
% Pages 9-20: a typescript, « Le théorème de Griffiths par voie algébrique »
% (title in pencil, underlined), paginated by hand 2-12 from p. 10 (p. 9, its
% first leaf, unnumbered), breaking off in the middle of typed p. 12 with the
% proof of Cor. 7.3. DECISION: transcribed in full as an unpublished typescript
% of his. It is none of his published SGA 7 exposés (I, II, VI-IX): the SGA 7
% plan of folder 35 has « XX Th. de Griffiths » struck out, and none of those
% exposés treats Griffiths' theorem. Evidence that it is his: written in the
% first person by an author who is neither Deligne nor Katz (« Deligne a
% vérifié... », « C'est ce que Katz vérifie... », « 7. Le procédé de Katz. A lui
% de jouer .... », then « Je vais prouver la conjecture du n° 6 ... en adaptant
% la méthode utilisée par Katz »); it speaks of « la théorie des cycles
% évanescents, qui sera prouvée dans SGA 7 cette année », of the « conjectures
% standart » and the « yoga des motifs », of the Lefschetz theorem « vache »
% (his term elsewhere in the fonds: folders 14, 23, 44, 162-1), « en l'état
% actuel de la Science »; and it sits in the folder he titled « SGA 7 (ma
% part) ». The ink corrections and margin notes (« non, pas tout à fait ! »,
% « pas prouvé », the boxed note on U' and Bertini) are taken as his, in the
% same pen as the \add{} revisions; typed interlinear insertions are given plain
% and noted when an ink stroke calls them in; machine x-outs are not recorded
% unless they read. The machine's ell is a chi-shaped glyph, set \ell.
%
% Pagination and numbered runs: remarks leaves 1-6 (pp. 2-7), sections
% (1)-(11) circled; typescript pp. [1]-12 (pp. 9-20), numbered sections 1-7 with
% Théorème 7.1, Corollaires 7.2, 7.3. Section 7 is written twice: the first
% « 7. Le procédé de Katz. A lui de jouer .... » closes typed p. 8 (p. 16),
% the second opens typed p. 9 (p. 17).
%
% What the batch is about. Typescript: 1, the class u(x) in H^1(S, R^{m-1}f_*)
% of a primitive class x, via Leray, and its compatibility with cohomological
% correspondences; applied to cycles Z on Y x_S C and z = Z(t_1) - Z(t_0);
% 2, for i : Y -> X over k algebraically closed, under a) (an algebraic
% right inverse to i^*, a case of the standard conjectures) and b) (no level-1
% piece in Coker i^*), every z algebraically equivalent to 0 decomposes as
% i^*(z') + z''; 3, the relative version; 4, the Griffiths homomorphism
% P^m(X) -> H^1(S, R^{m-1}f_*Q_ell) for a Lefschetz pencil; 5, hence the
% Griffiths group of a generic hyperplane section contains P^{2n}_alg(X)
% (tensor Q); 6, conjecture: for hypersurface sections of large degree,
% condition b) holds — proved in char. 0 by Lefschetz-Griffiths via
% H^{2n-1}(O), in char. p reduced to finding a Frobenius eigenvalue that is a
% p-adic unit; 7, Th. 7.1 (Galois action on the Frobenius-fixed points of the
% semisimple part of coker H^{n-1}(X,O) -> H^{n-1}(Y,O) has elements of every
% trace a), Cor. 7.2 (card Y_t(k') = a mod p via Deligne's congruence and
% Chebotarev), Cor. 7.3 (vanishing cycles not of level 1 on generic lines).
% The ink marks 7.1 and 7.2 « pas prouvé ».
%
% Notation fixed once: \ell for the typed chi; \underline{\mathbf{Z}}_\ell,
% \underline{\mathbf{Q}}_\ell, \underline{\mathcal{O}} as the machine
% underlines; E_0(Y), V, V_0, G_0 as on p. 17. Readings to check first: the
% top-margin note of p. 17 (« non »), the two slanting margin notes of p. 17,
% the boxed margin note of p. 19, the inked formula heading p. 20, and the
% word « pour » over a remark on p. 3.

\documentclass[11pt,a4paper]{article}
% Le préambule commun à toutes les transcriptions du fonds Grothendieck.
%
% Il tient en une idée : **l'incertitude se compose, elle ne se résout pas.**
% Une transcription qui lisse un mot douteux a détruit la seule chose pour
% laquelle on la faisait — un lecteur doit pouvoir savoir, à chaque endroit,
% ce qui était sur le feuillet et ce qui a été ajouté. Les six macros qui
% suivent sont donc l'appareil critique, et non de la décoration.
%
% Les mêmes commandes sont reconnues par `scripts/render.mjs`, qui produit la
% vue de lecture du site. Une macro ajoutée ici doit l'être aussi là-bas,
% faute de quoi le rendu échouera bruyamment — ce qui est voulu : un rendu
% silencieusement faux est pire qu'un rendu absent.

\ProvidesPackage{grothendieck}[2026/08/08 Transcriptions du fonds Grothendieck]

\usepackage{amsmath,amssymb,amsthm}
\usepackage{mathrsfs}
\usepackage{stmaryrd}
% Les diagrammes commutatifs sont partout dans ce fonds : c'est la langue
% même de la géométrie algébrique de Grothendieck, et une transcription qui
% les rendrait en prose ne transcrirait rien.
\usepackage{tikz-cd}
% Français par défaut — c'est la langue des feuillets — mais l'anglais est
% chargé aussi : la traduction et le résumé sont des documents anglais, et la
% typographie française (espaces avant les deux-points) y serait une faute.
% Ces documents basculent par \selectlanguage{english} en tête de corps.
\usepackage[english,french]{babel}
\usepackage{fontspec}
\usepackage{geometry}
\geometry{a4paper,margin=25mm,marginparwidth=28mm}
\usepackage{marginnote}
\usepackage{xcolor}
% ulem, not soul. `soul`'s \st and \ul are built on a character-by-character
% scan that XeLaTeX's font machinery breaks: the document fails with an
% undefined control sequence rather than degrading. ulem does the same two jobs
% and is Unicode-engine safe. `normalem` stops it hijacking \emph, which the
% transcriptions use for Grothendieck's own emphasis.
\usepackage[normalem]{ulem}
\usepackage{hyperref}
\hypersetup{colorlinks=true,linkcolor=black,urlcolor=black,pdfborder={0 0 0}}

% --- Métadonnées du lot -----------------------------------------------------
% Reprises telles quelles par le script de rendu, qui les affiche en tête de
% la vue de lecture. Elles fixent ce que le fichier prétend couvrir : sans
% elles, une transcription flotte sans point d'attache dans un fonds de seize
% mille feuillets.

\newcommand{\folder}[1]{\gdef\grothFolder{#1}}
\newcommand{\batch}[1]{\gdef\grothBatch{#1}}
\newcommand{\leaves}[2]{\gdef\grothFirst{#1}\gdef\grothLast{#2}}
\newcommand{\foldertitle}[1]{\gdef\grothFoldertitle{#1}}
\gdef\grothFolder{?}\gdef\grothBatch{?}\gdef\grothFirst{?}\gdef\grothLast{?}\gdef\grothFoldertitle{}

% --- Le résumé --------------------------------------------------------------
%
% Il ouvre la lecture modernisée, sous le titre « Résumé », et il est composé
% en petit. Ce n'est pas un ornement : c'est le seul passage du document qui
% ne soit pas des mathématiques, et un lecteur doit pouvoir le reconnaître
% avant de l'avoir lu — puis le sauter s'il connaît déjà le sujet, ou s'y
% arrêter s'il ne le connaît pas. Le filet qui le suit marque la reprise du
% corps.

\newenvironment{resume}
  {\small\setlength{\parskip}{0.45em}}
  {\par\medskip\hrule\medskip}

% --- Mots-clés ---------------------------------------------------------------
%
% En anglais, en clôture du résumé de la lecture modernisée : le vocabulaire
% moderne sous lequel chercher ce que les feuillets construisent. Le manifeste
% du site (`npm run manifest`) les extrait pour étiqueter la cote entière —
% cette ligne est la source unique des tags, il n'y a pas de fichier à côté.
\newcommand{\keywords}[1]{\par\smallskip\noindent
  {\footnotesize\textsc{Keywords} --- \textit{#1}}\par}

% --- L'appareil critique ----------------------------------------------------

% Le numéro de feuillet, en marge. C'est l'ancre qui permet de régler un
% désaccord sur une formule en regardant *un* feuillet plutôt qu'une chemise
% de six cents. Le site s'en sert aussi pour tourner les pages du fac-similé
% quand on fait défiler la transcription.
\newcommand{\leaf}[1]{%
  \marginnote{\footnotesize\color{gray!70}\textsf{#1}}%
  \label{leaf:#1}\ignorespaces}

% Illisible : on ne devine pas. Un mot inventé qui se lit comme les autres est
% le pire résultat possible.
\newcommand{\ill}{\textcolor{red!60!black}{[\,\ldots\,]}}

% Lecture douteuse : le mot est donné, et signalé comme incertain.
\newcommand{\uncertain}[1]{\uline{#1}}

% Ajout de l'éditeur — une lettre manquante, un mot sous-entendu.
\newcommand{\add}[1]{\textcolor{blue!50!black}{[#1]}}

% Biffé par Grothendieck lui-même. Ce qu'il a rayé fait partie du document :
% une rature dit souvent où la pensée a changé de direction.
\newcommand{\struck}[1]{\sout{#1}}

% Note du transcripteur, sur l'état matériel du feuillet.
\newcommand{\note}[1]{\par\smallskip{\small\color{gray!60!black}\textit{[#1]}}\par\smallskip}

% Note marginale de Grothendieck, distincte du corps du texte.
\newcommand{\marginal}[1]{\par\smallskip\noindent\hspace{1em}%
  {\small\color{gray!60!black}#1}\par\smallskip}

% --- Titre ------------------------------------------------------------------

\AtBeginDocument{%
  \begin{center}
    {\large\bfseries Fonds Alexandre Grothendieck --- cote n\textsuperscript{o}~\grothFolder}\\[2pt]
    {\itshape\grothFoldertitle}\\[4pt]
    {\small feuillets \grothFirst--\grothLast\ (lot \grothBatch)}\\[2pt]
    {\footnotesize\color{gray!60!black}Transcription établie d'après le fac-similé numérisé
     par l'Université de Montpellier.\\
     Les lectures incertaines sont soulignées, les illisibles notées [\,\ldots\,],
     les ajouts entre crochets.}
  \end{center}
  \vspace{1em}}

\endinput


\folder{36}
\batch{1}
\pages{1}{20}
\dating{[vers 1967-1973]}
\watermark{Édition de démonstration}
\foldertitle{SGA 7 (ma part) : notes manuscrites (s.d.), tapuscrit annoté (s.d.)}

\begin{document}

\page{2}
\note{feuillet 1 sur 6 (numéroté « 1 » en haut à droite) de remarques d'une
autre main, non identifiée, à l'encre bleue sur papier quadrillé, titrées
« Remarques. SGA 7 IX » : un lecteur relit un état dactylographié de l'exposé IX
de SGA 7 (modèles de Néron et monodromie) et renvoie, par page et ligne, à ce
tapuscrit, en tutoyant son auteur. Elles ne sont pas reproduites. Ce feuillet
couvre l'introduction (numéros 0 et i à iv) et le §1, pages 0 à 2 du tapuscrit :
$\ell$ premier à $p$, pourquoi $S$ hensélien, « à isogénie près », réduction
semi-stable d'une courbe, la définition d'isogénie, $A \times_S
\operatorname{Spec}(K) \simeq A_K$. Réponses de l'auteur : aucune en toutes
lettres ; au crayon, des cercles autour de plusieurs renvois de ligne, et un
ovale suivi d'un point d'interrogation autour de « lié à la monodromie ». Ce
crayon est vraisemblablement de sa main, sans que cela soit établi}

\page{3}
\note{feuillet 2 sur 6 (« 2 ») des mêmes remarques, non reproduites : fin du
§1 et §2 (pages 3 à 20 du tapuscrit) — l'accouplement d'A.~M.\ comparé à celui
de l'auteur, « formellement non ramifié », « séparablement clos », un argument
de la page 14 jugé « un peu immoral » (un morphisme entre extensions d'une
variété abélienne par un tore est strictement compatible à la filtration par le
poids, et est une isogénie si et seulement si $T_\ell(\varphi)$ en est une),
Rmq 2.6 a) et le passage à la limite pour les modèles de Néron. Réponses de
l'auteur : aucune en toutes lettres ; au crayon, des cercles autour de « pg 20 »
et de renvois de ligne, et au-dessus d'un renvoi de la page 8 un mot entouré,
\uncertain{pour}, suivi d'un point d'interrogation. Main vraisemblablement la
sienne, non établie}

\page{4}
\note{feuillet 3 sur 6 (« 3 ») des mêmes remarques, non reproduites : §3 et §4
(pages 23 à 40 du tapuscrit) — ${}_nK$ étale séparé comme noyau de ${}_nG \to
{}_nH$, les notions de 3.5 pour $S$ non hensélien, « clôture normale », un énoncé
de « voisinage tubulaire » en 4.2 (morphisme de topos
$\operatorname{Spec}(K)_{\mathrm{et}} \to \operatorname{Spec}(k)_{\mathrm{et}}$ et
$A^0_0$ schéma relatif annelé par $\varphi^{*}\mathcal{O}$), la trace
contragrédiente en page 36, le sens de « $\ell$-groupe » en 4.5.4. Réponses de
l'auteur, au crayon : en marge de la remarque sur la page 33, un grand point
d'interrogation entouré ; en marge de la remarque sur la page 36 (« ce n'est pas
la trace qui est contragrédiente »), entouré : « si ? » ; des cercles autour de
« pg 28 », « dim » et « pas ». Main vraisemblablement la sienne, non établie}

\page{5}
\note{feuillet 4 sur 6 (« 4 ») des mêmes remarques, non reproduites : fin du §4
et §§5--6 (pages 41 à 63 du tapuscrit) — la « fonction de Dirac » sur un
groupe, 4.7 en dehors du cas strictement local, le sens de (5.2.3.1), 5.7, 5.11,
5.15, l'usage de $\operatorname{car}(K)=0$ via Tate, $A^0$ $\ell$-divisible ou
non, et une objection à la page 59 : la plus grande sous-variété abélienne
ordinaire ne se relèverait pas, et la page 63 serait fausse (déformer $E \times
F$, $E$ ordinaire, $F$ non). Réponses de l'auteur : en marge de la page 41, au
crayon : « guillemets suffisants ? » ; en marge de la remarque sur la page 46,
au crayon et entouré : « je le dis ! » ; en regard de « pg 56 bis 5.15 », entouré
au crayon, un grand point d'interrogation ; en marge de la remarque sur la page
55 ter, au crayon et entouré : « c'est dit dans la 5.13 » ; en marge de
l'objection de la page 59, à l'encre : « CANULAR ?! ». Main vraisemblablement la
sienne, non établie}

\page{6}
\note{feuillet 5 sur 6 (« 5 ») des mêmes remarques, non reproduites : fin du §6
et §§7, 9, 10 (pages 60 à 98 du tapuscrit ; le numéro 8 est laissé vide) — le
« il est immédiat » de 6.4, l'existence à isogénie près de « parties finies » et
« parties toriques » du bon rang, la référence à Serre pour l'indépendance de
$\ell$, $A_0$ non défini page 65, 9.2.3, 9.3.3, $S$ strictement hensélien en 9.4
et page 92 bis, un signe en 9.5.8, 10.5.3--10.5.4. Réponses de l'auteur, au
crayon et entourées : dans un ovale qui prend « pg 60 » et « pg 61, 61 bis : je
n'ai pas lu » : « c'est pourtant trivial » ; en regard des « parties finies » :
« non, mon argument marche » ; en regard de la remarque sur la page 66 : « que
faire ? » ; en regard de la page 67, un point d'interrogation. Main
vraisemblablement la sienne, non établie}

\page{7}
\note{feuillet 6 sur 6 (« 6 ») des mêmes remarques, non reproduites : §11
(pages 101 et 103 du tapuscrit) — où trouver que le modèle de Néron commute à
$S \mapsto \hat{S}$, et le prolongement d'une biextension de $A \times A'$ ramené
à celui des ${}_{\ell^\infty}A \times {}_{\ell^\infty}A'$. Réponse de l'auteur :
en marge de la remarque sur la page 103 (« le renvoi au courage du lecteur est
désagréable »), au crayon, entouré avec « pg 103 » : « que faire ? ». Main
vraisemblablement la sienne, non établie. Le reste du feuillet est vierge}


\section*{Le théorème de Griffiths par voie algébrique}

\page{9}

\note{titre écrit au crayon en tête du premier feuillet du tapuscrit, et souligné
; un grand crochet au crayon brun le précède dans la marge. Le tapuscrit est
paginé à la main en haut à droite, de 2 (page 10) à 12 (page 20) ; ce feuillet-ci,
le premier, ne porte pas de numéro. Le $\ell$ de la machine, partout, est un
caractère en forme de $\chi$ ; les faisceaux constants $\underline{\mathbf{Z}}_\ell$,
$\underline{\mathbf{Q}}_\ell$ sont soulignés à la machine. Les ajouts à la main
sont donnés par \add{}, avec une note}

1.\quad Soit
\[ f\colon Y \longrightarrow S \]
un morphisme propre et lisse de schémas, avec $S$ \add{régulier} connexe à
caractéristiques résiduelles toutes $\neq \ell$ (nombre premier fixé).\note{« régulier
» ajouté à l'encre au-dessus de la ligne} Soit $s \in S$, $\bar{s}$ un point
géométrique au dessus de $s$. Soit $F$ un faisceau $\ell$-adique constant tordu sur
$S$ (en pratique, un faisceau de la forme $\underline{\mathbf{Z}}_\ell(i)$), et $x$
une classe de cohomologie
\[ x \in H^m(Y, F_Y) . \]
Disons que $x$ est \emph{primitive} si son image dans $H^0(S, R^m f_{*}(F))$ est
nulle, ou ce qui revient au même par le th.\ de spécialisation de la cohomologie,
si son image dans $H^m(Y_{\bar s}, F_{\bar s})$ est nulle.\note{dans ces trois
formules l'exposant $m$ est écrit ou repassé à la main sur une autre lettre
dactylographiée (un $N$, semble-t-il), et une parenthèse dactylographiée à la
suite, « $(F)$ », est noircie ; de même $m-1$ ci-dessous, et $f_{*}(F_Y)$ y est
repassé à l'encre} Alors la suite spectrale de Leray pour $f$ nous donne un
élément
\[ u(x) \in E_2^{1,m-1} = H^1(S, R^{m-1} f_{*}(F_Y)) . \]
Donnons une propriété fonctorielle générale de cette classe $u(x)$. Pour ceci, soit
$Y'$ lisse et propre sur $S$, et
\[ Z \in H^N(Y \times_S Y', G), \]
de sorte que par le formalisme des correspondances cohomologiques, $Z$ définit une
application
\[ x' \longmapsto Z(x')\colon H^{m'}(Y', F') \longrightarrow H^m(Y, F), \]
si on suppose $m = m' + N - \struck{\dim(Y'/S)}$, et qu'on a un
accouplement $F' \otimes G \to F(d')$\note{après « $N-$ », la machine a biffé
« $\dim(Y'/S)$ » sans rien mettre à la place ; le « $d'$ » que la suite définit
manque dans la formule} (NB $d'$ est la dimension relative de $Y'$ sur $S$,
supposée constante, et $(d')$ désigne le twist de Tate ; $F'$ et $G$ désignent des
faisceaux $\ell$-adiques constants tordus \add{sur $S$} tout comme $F$).\note{« sur
$S$ » ajouté à l'encre} Ceci dit, si $x'$ est primitive, alors il en est
évidemment de même de $x = Z(x')$, et on aura de plus
\[ u(Z(x')) = Z(u(x')) , \]
où le deuxième membre s'explicite en notant que $Z$ définit également, par

\page{10}

\note{p.~2 du tapuscrit}

localisation étale sur $S$, un homomorphisme
\[ a' \longmapsto Z(a')\ :\ R^{m'-1} f'_{*}(F_{Y'}) \longrightarrow R^{m-1}
f_{*}(F_Y) , \]
et par suite un homomorphisme (encore noté par $Z$) des $H^1(S,\ )$ pour les deux
membres.

Nous allons appliquer ceci au cas où $Y'$ est une courbe relative $C$ sur $S$,
$m'=2$, $x'$ est la classe de cohomologie
\[ x' \in H^2(C, \underline{\mathbf{Z}}_\ell(1)) \]
défini par un diviseur relatif de la forme $t_1(S) - t_0(S)$, où $t_1$ et $t_0$
sont deux sections de $C$ sur $S$. Supposant $C$ à fibres géométriques connexes on
voit donc que $x'$ est primitif. (NB on pourrait prendre aussi bien n'importe quel
diviseur relatif de degré relatif nul.) On prendra pour $Z$ la classe de
cohomologie
\[ Z \in H^{2n}(Y \times_S C, \underline{\mathbf{Z}}_\ell(n)) \]
associée à un cycle algébrique sur $Y \times_S C$ de codimension $n$, qu'on notera
aussi $Z$. (NB pour que cette classe de cohomologie soit vraiment définie en l'état
actuel de la Science, il faut faire des hypothèses restrictives genre $S$ de type
fini sur un corps, ou $S$ excellent de caractéristique nulle ou le cycle $Z$ un
cycle « relatif » i.e.\ de codimension $\leqslant n$ fibre par fibre. Sous ces
conditions, $x$ est la classe de cohomologie
\[ x \in H^{2n}(Y, \underline{\mathbf{Z}}_\ell(n)) \]
associée au cycle algébrique
\[ z = Z(t_1) - Z(t_0) \]
sur $Y$. (NB on suppose que les intersections écrites ont un sens, ou qu'on
travaille avec des \emph{classes} de cycles de façon qu'elles en aient un, -- ce
dernier cas supposant (toujours en l'état actuel de la Science) que $S$ est de type
fini sur un corps.

\page{11}

\note{p.~3 du tapuscrit}

Ceci posé, le cycle $Z$ (ou plutôt, sa classe de cohomologie) définit un
homomorphisme naturel
\[ (*) \qquad Z\colon\ R^1 g_{*}(\underline{\mathbf{Z}}_\ell(1)) \longrightarrow
R^{2n-1} f_{*}(\underline{\mathbf{Z}}_\ell(n)) , \]
d'où un homomorphisme
\[ (**) \qquad Z\colon\ H^1(S, R^1 g_{*}(\underline{\mathbf{Z}}_\ell(1)))
\longrightarrow H^1(S, R^{2n-1} f_{*}(\underline{\mathbf{Z}}_\ell(n))) , \]
et la compatibilité générale énoncée plus haut nous dit que $u(x) =
u(c\ell(z))$ n'est autre que l'image de $u(x')$ du premier membre de $(**)$.\note{$g$
est le morphisme structural de $C$, que le tapuscrit n'a pas nommé} Nous aurons
besoin simplement pour la suite de savoir que \emph{si l'homomorphisme $(*)$ est
nul (ou ce qui revient au même, si l'homomorphisme correspondant sur les fibres
géométriques en $\bar s$}
\[ (*\ \mathrm{bis}) \qquad H^1(C_{\bar s}, \underline{\mathbf{Z}}_\ell(1))
\longrightarrow H^{2n-1}(Y_{\bar s}, \underline{\mathbf{Z}}_\ell(n)) \]
\emph{est nulle), alors il en est de même de} $u(x) \in H^1(S, R^{2n-1}
f_{*}(\underline{\mathbf{Z}}_\ell(n)))$.\note{les passages en italique sont
soulignés, à la main semble-t-il}

2.\quad Soit
\[ i\colon Y \longrightarrow X \]
un morphisme de schémas propres et lisses sur un corps algébriquement clos $k$, $n$
un entier.\note{« propres » est tapé sur un autre mot, biffé à la machine} On
\struck{suppose} considère l'homomorphisme
\[ (***) \qquad i^{*}\colon\ H^{2n-1}(X, \underline{\mathbf{Q}}_\ell(n))
\longrightarrow H^{2n-1}(Y, \underline{\mathbf{Q}}_\ell(n)) , \]
et on suppose :\note{« $({*}{*}{*})$ » est écrit à la main dans la marge ; dans les
deux membres un $\mathbf{Q}$ est tapé sur le $\mathbf{Z}$}

a)\quad Il existe un homomorphisme en sens inverse $v$, dont la restriction à
l'image de $i^{*}$ soit un inverse à droite de $i^{*}$, et qui soit « algébrique »
i.e.\ donné par un cycle algébrique sur $Y \times X$. (NB Cette hypothèse est un
cas particulier des « conjectures standart », donc est sans doute toujours
vérifiée ; elle est vérifiée en tous cas si $i$ est l'inclusion d'une section
hyperplane dans $X$, lorsque l'opération $\Lambda$ sur la cohomologie de $X$ est
algébrique, \add{propriété} qu'on sait vérifier dans de très nombreux cas, par
exemple lorsque $X$ se remonte en car.\ zéro avec sa polarisation.)\note{« propriété
» est tapé dans l'interligne, au-dessus d'un mot biffé à la machine ; de même, un
« $X$ » biffé à la machine suit « sur »}

\page{12}

\note{p.~4 du tapuscrit}

b)\quad Pour toute courbe algébrique (propre, lisse, connexe) $C$ sur $k$, et tout
cycle algébrique $Z$ sur $Y \times C$ de codimension $n$, l'homomorphisme
correspondant
\[ Z\colon\ H^1(C, \underline{\mathbf{Q}}_\ell(1)) \longrightarrow H^{2n-1}(Y,
\underline{\mathbf{Q}}_\ell(n)) \]
a son image dans celle de $i^{*}$ de $(***)$. (NB Si on admet le yoga des motifs,
justifiable par les « conjectures standard », cette condition signifie que le
\add{plus grand} morceau « de niveau 1 » du motif ayant comme réalisation
$\ell$-adique $\operatorname{Coker} i^{*}$ est nul).\note{« plus grand » ajouté à
l'encre au-dessus de « morceau »}

Sous ces conditions, \emph{pour tout cycle algébrique $z \in Z^n(Y)$ qui est
algébriquement équivalent à zéro, il existe un cycle $z' \in Z^n(X)$
algébriquement équivalent à zéro, et un $z'' \in Z^n(Y)$, tels que}
\begin{itemize}
\item[a')] $z = i^{*}(z') + z''$
\item[b')] $z'' = Z''(t_1) - Z''(t_0)$, où $Z''$ est un \supplied{$\mathbf{Q}$-}cycle
\emph{de codimension $n$ sur un produit $Y \times C$, induisant un homomorphisme
nul} $H^1(C, \underline{\mathbf{Q}}_\ell(1)) \to H^{2n-1}(Y,
\underline{\mathbf{Q}}_\ell(n))$.
\end{itemize}
\note{« $\mathbf{Q}$- » ajouté à l'encre au-dessus de « cycle » ; dans
« $Y \times C$ », le $Y$ est écrit à l'encre sur un $X$ biffé ; les soulignements
sont à la main}

Démonstration immédiate : on écrit $z = Z(t_1) - Z(t_0)$ avec $Z \in Z^n(Y \times
C)$, ce $Z$ définit un homomorphisme $Z\colon H^1(C) \to H^{2n-1}(Y)$, qui par
hypothèse b) tombe dans $\operatorname{Im} i^{*}$ donc par hypothèse a) est égal
au composé $(iv)Z = i(vZ)$, or $vZ\colon H^1(C) \to H^{2n-1}(X)$ est algébrique
puisque $v$ l'est, donc donné par un cycle algébrique $Z'$ sur $X \times C$, et si
on pose $Z'' = Z - (i \times \mathrm{id}_C)^{*}(Z')$, on aura par construction que
$Z''$ induit l'homomorphisme nul $H^1(C) \to H^{2n-1}(Y)$. Il suffit alors de
poser
\[ z' = Z'(t_1) - Z'(t_0), \qquad z'' = Z''(t_1) - Z''(t_0) . \]
\note{dans « $(iv)Z = i(vZ)$ » le tapuscrit écrit $i$ pour $i^{*}$ ; on le laisse}

3.\quad Mettant ensemble les réflexions de 1) et 2), on trouve ceci. On se donne
des schémas propres et lisses $X$ et $Y$ sur $S$ \add{régulier} connexe, et un
$S$-morphisme
\[ i\colon Y \longrightarrow X . \]
\note{« régulier » ajouté à l'encre, comme page 9}

\page{13}

\note{p.~5 du tapuscrit. La première ligne est tapée dans l'interligne et appelée
par un trait à l'encre à la place d'un passage biffé à la machine}

On suppose \add{a) et b) vérifiés pour le morphisme $i_{\bar s}\colon X_{\bar s}
\to X_{\bar s}$ induit} pour les fibres en $\bar s$ (point géométrique de $S$ au
dessus du point \emph{générique} de $S$).\note{ainsi, $X_{\bar s} \to X_{\bar s}$ :
on attendrait $Y_{\bar s} \to X_{\bar s}$. Deux lignes suivantes, entièrement
biffées à la machine, ne se lisent pas (\struck{\ill{}}) ; « Soit » est tapé
au-dessus de leur fin} Soit $z \in Z^n(Y)$ tel que le cycle induit sur $Y_{\bar s}$
soit algébriquement équivalent à zéro, et à fortiori cohomologiquement équivalent
à zéro. Donc la classe
\[ u(z) \in H^1(S, R^{2n-1} f_{*}(\underline{\mathbf{Z}}_\ell(n))) \]
est définie. Ceci posé, je dis que \emph{cette classe est contenue dans}
$i^{*}(H^1(S, R^{2n-1} g_{*}(\underline{\mathbf{Z}}_\ell(n))))$ \emph{modulo une
classe de torsion}, du moins si on se permet de remplacer $S$ par un ouvert non
vide (restriction \add{probablement} inutile, mais cet énoncé nous suffira pour ce
qu'on veut en faire).\note{« cette classe est » est suivi de mots biffés à la
machine (« une classe de torsion », semble-t-il) ; « probablement » est tapé
au-dessus d'un mot biffé ; les soulignements sont à la main. Sous la formule, une
ligne biffée à la machine ; une autre, biffée aussi, avant « se permet »}

Pour le voir, on applique 2) au morphisme $i_{\bar s}$. La courbe $C$ et le cycle
$Z''$ peuvent se définir déjà sur une extension finie $K$ de $k(s)$, Un argument
de trace permet de supposer que $K = k(s)$, et \struck{\ill{}} on peut supposer que
$C$ est une courbe propre et lisse relative sur $S$ lui-même, et $Z''$ un cycle sur
$Y \times_S C$, et que la formule $z = i^{*}(z') + z''$ est valable sur $Y \times_S
C$. Elle implique une décomposition analogue de $u(z)$, qui est la somme de
$i^{*}(u(z'))$ et de $u(z'')$, de sorte qu'il suffit de prouver que $u(z'')$ est
nul. Or $u(z'')$ induit par construction 2) l'homomorphisme nul $H^1(C_{\bar s})
\to H^{2n-1}(Y_{\bar s})$. Il résulte donc bien de 1) que $u(z'')$ est nul, cqfd.

4.\quad Soient $X$ un schéma projectif et lisse de dimension $m$ \add{sur $k$ alg.\
clos}, prenons un pinceau de Lefschetz, de façon à obtenir une variété éclatée non
moins projective et lisse $\widetilde{X}$ et un morphisme\note{« sur $k$ alg.\ clos
» ajouté à l'encre au-dessus de la ligne}
\[ \tilde f\colon \widetilde{X} \longrightarrow \mathbf{P}^1_k \]
dont les fibres sont les sections hyperplanes du pinceau. Soit $S$ un ouvert non
vide de $\mathbf{P}^1_k$ au dessus duquel $\widetilde{X}$ soit lisse, et soit
$f\colon Y = \widetilde{X}|S \to S$\note{fin de ligne en partie repassée à la
main ; la lettre de l'ouvert est tapée sur une autre, et on lit $S$ sous réserve}

\page{14}

\note{p.~6 du tapuscrit}

Soit $P^m(X)$ la partie primitive de la cohomologie $H^m(X,
\underline{\mathbf{Q}}_\ell)$, i.e.\ le noyau du composé
\[ H^m(X) \to H^m(\widetilde{X}) \longrightarrow H^m(Y) \longrightarrow H^0(S, R^m
f_{*}(\underline{\mathbf{Q}}_\ell)) . \]
On trouve donc grâce à 1) un homomorphisme (dit « de Griffiths »)
\[ P^m(X) \longrightarrow H^1(S, R^{m-1} f_{*}(\underline{\mathbf{Q}}_\ell)) \]
(en fait, on y pourrait remplacer $S$ par $\mathbf{P}^1_k$, mais peu importe ici).
D'ailleurs on a un homomorphisme d'immersion canonique

\begin{tikzcd}
Y \arrow[rr, "i"] \arrow[dr, "f"'] & & X_S \arrow[dl, "g"] \\
& S &
\end{tikzcd}

induisant
\[ i^{*}\colon\ H^1(S, R^{m-1} g_{*}(\underline{\mathbf{Q}}_\ell)) = H^1(S,
H^{m-1}(X, \underline{\mathbf{Q}}_\ell)) \longrightarrow H^1(S, R^{m-1}
f_{*}(\underline{\mathbf{Q}}_\ell)) . \]
Ceci dit, \emph{l'image inverse par l'homomorphisme de Griffiths plus haut de $\operatorname{Im} i^{*}$ est réduite à zéro}.

Ceci se fait par des calculs explicites essentiellement triviaux, compte tenu de la
structure cohomologique connue d'une variété éclatée.

5.\quad Mettant ensemble 3) et 4), on trouve ceci :

\struck{\ill{}} Avec les notations de 4), \add{supposant $m = 2n$,} soit $\bar s$ le
point générique de $\mathbf{P}^1_k$, et \emph{supposons que l'immersion $Y_{\bar
s} \to X \otimes_k k(\bar s)$ de la section hyperplane générique satisfasse aux
conditions a) et b) de 2).} \struck{\ill{}} \emph{Alors pour tout $z \in Z^n(X)$
dont la restriction à $Y_{\bar s}$ est algébriquement équivalent à zéro, $z$ est
cohomologiquement équivalent à zéro} (mod torsion).\note{« supposant $m=2n$, » est
tapé dans l'interligne et appelé par un trait à la main, à la place d'une ligne
biffée à la machine (« si on suppose de plus $m = 2n$ », semble-t-il) ; les
soulignements sont à la main} \struck{\ill{}} Par suite, le groupe quotient du
groupe des cycles cohomologiquement équivalents à zéro \struck{\ill{}} sur $Y_{\bar
s}$, modulo ceux qui sont algébriquement équivalents à zéro (groupe qui mérite le
nom de \emph{groupe de Griffiths} $\operatorname{Grif}(Y_{\bar s})$ de $Y_{\bar
s}$), tensorisé par $\underline{\mathbf{Q}}$, contient comme sous-espace le
sous-espace $P^{2n}_{\mathrm{alg}}(X)$ formé des classes de cohomologie primitives
sur $X$ qui sont \emph{algébriques}, -- vectoriel qui a

\page{15}

\note{p.~7 du tapuscrit}

une nette tendance à être $\neq 0$.

6.\quad \struck{Conjecture II} \emph{Conjecture} : pour $X$ \add{$\subset
\mathbf{P}^r_k$} fixé comme dans 5), \add{avec $n \geqslant 2$,} il existe un
entier $d_0$ tel que pour tout pinceau \add{de Lefschetz} \add{assez général}
\struck{de} hypersurfaces de degré $d \geqslant d_0$ dans $\mathbf{P}^r$, donnant
une section générique $Y_{\bar s}$ avec $X$, l'inclusion $Y_{\bar s}
\hookrightarrow X_{\bar s}$ satisfait à la condition b) de 2) (le fait que la
condition a) doive être vérifiée en tous cas étant de toutes façons justiciable de
la conjecture d'algébraicité de la correspondance cohomologique $\Lambda$ sur
$X$).\note{« $\subset \mathbf{P}^r_k$ », « avec $n \geqslant 2$ » et « de
Lefschetz » sont tapés dans l'interligne et appelés à l'encre ; « assez général »
est à l'encre ; le « 5 » de « 5) » est tapé sur un autre chiffre. Un double trait
vertical, à l'encre, marque en marge l'énoncé de la conjecture}

Nous allons examiner cette conjecture, en supposant déjà que le théorème de
Lefschetz « vache » soit vrai sur $X$, et que la correspondance $\Lambda$ sur $X$
est algébrique. La théorie des cycles évanescents, qui sera prouvée dans SGA 7
cette année, donne alors ceci : L'orthogonal $E^{2n-1}$ de $H^{2n-1}(X)$, considéré
comme plongé dans $H^{2n-1}(Y_{\bar s})$, est un \emph{module simple sous l'action
de} $\operatorname{Gal}(k(\bar s)/k(s))$, ainsi que \add{sous} celle de tout
sous-groupe \struck{d'indice fini} \add{ouvert}.\note{« sous » et « ouvert » à
l'encre, le second au-dessus de « d'indice fini » barré} (NB pour ce dernier fait,
on utilise de façon essentielle qu'on est en dimension impaire $2n-1$, de sorte que
les opérations de monodromie locale sont des transvections et non des symétries,
donc ne se tuent pas en passant à une puissance quelconque \ldots) Il s'ensuit déjà
que la condition b) \add{équivaut à dire} que $E^{2n-1}$ \emph{n'est pas tout
entier de niveau 1}. En caractéristique nulle, s'il était tout entier de niveau 1,
\add{cela impliquerait} que les composantes de type $p,q$ de $E^{n-1}$ sont nulles
sauf les deux composantes $(n,n-1)$ et $(n-1,n)$ (la réciproque étant vraie
d'ailleurs si la conjecture de Hodge est vraie) ; cela signifie aussi que
$H^{p,q}(X) \hookrightarrow H^{p,q}(Y)$ est un isomorphisme pour ces valeurs de
$(p,q)$, en particulier que $H^{2n-1}(X, \underline{\mathcal{O}}_X) \to
H^{2n-1}(Y, \underline{\mathcal{O}}_Y)$ est un isomorphisme \add{(car $n
\geqslant 2$)}, ce qui n'est pas le cas pour $d$ grand (car la dimension du second
membre tend vers $\infty$ avec $d$). Cela prouve (par la méthode de Lefschetz,
reprise par Griffiths) l'énoncé voulu en caractéristique nulle.\note{« équivaut à
dire » et « cela impliquerait » sont tapés dans l'interligne au-dessus de passages
biffés à la machine ; d'autres biffures à la machine, illisibles, ne sont pas
relevées ; « $(\text{car } n \geqslant 2)$ » est à l'encre. « $E^{n-1}$ » est bien
ce que porte le tapuscrit, au lieu de $E^{2n-1}$}

\page{16}

\note{p.~8 du tapuscrit}

En caractéristique quelconque, on ne peut pour l'instant que donner des cas assez
particuliers où la conjecture soit vérifiée ; par Katz, cela marche en tous cas pour
\struck{les} \add{$X$} intersection\struck{s} complète\struck{s} \add{et un pinceau
« assez général »} (ce qui suffit pour faire marcher l'exemple minimal de
Griffiths : $X$ quadrique dans $\mathbf{P}^5$ \ldots).\note{« $X$ » au-dessus de
« les » barré, et « et un pinceau assez général », dans l'interligne,
sont à l'encre ; le pluriel des deux mots suivants est corrigé à la main} Pour
vérifier que $E^{2n-1}(Y_{\bar s})$ n'est pas de niveau 1, on constate aisément
qu'il suffit de vérifier qu'il existe \emph{un} point $t$ du pinceau tel que
$E^{2n-1}(Y_{\bar t})$ ne soit pas de niveau 1 (il faut utiliser pour ceci le
critère de Néron-Ogg-Chafarévitch de bonne réduction des VA ; la triste vérité est
que je ne sais pas \add{montrer de façon générale}, même en admettant les
conjectures standard, que le niveau d'un motif ne peut que diminuer par
spécialisation \ldots) Quitte à écrire $k$ comme limite d'anneaux de type fini au
sens absolu, la même observation est d'ailleurs valable, en se spécialisant sur un
corps fini. Or dans ce cas, le niveau 1 implique \add{(et est peut-être équivalent
à)} la divisibilité \add{par $q^{n-1}$} des valeurs propres de frobenius opérant
sur $E^{2n-1}(Y_{\bar t})$ \struck{par $q^{n-1}$} ($q = \operatorname{card}
k(t)$), et il suffit donc de trouver $t$ tel que l'une des valeurs propres de
frobenius soit une unité $p$-adique. C'est ce que Katz vérifie dans divers cas par
une utilisation astucieuse de la formule de Lefschetz-Woodshole-Verdier.\note{«
montrer de façon générale » est tapé dans l'interligne au-dessus d'un mot biffé à
la machine ; « (et est peut-être équivalent à) » et « par $q^{n-1}$ » sont à
l'encre, et « par $q^{n-1}$ » tapé en fin de phrase est barré à l'encre}

7.\quad Le procédé de Katz. A lui de jouer \ldots.

\note{le reste du feuillet est vierge}

\page{17}

\note{p.~9 du tapuscrit : une seconde rédaction du n\textsuperscript{o}~7, qui
remplace « Le procédé de Katz. A lui de jouer \ldots » de la page 16}

7.\quad Je vais prouver la conjecture du n\textsuperscript{o}~6 pour des pinceaux
de Lefschetz assez généraux d'hypersurfaces de grand degré, en adaptant la méthode
utilisée par Katz pour les intersections complètes. Nous allons
prouver :\note{« vais prouver » est entouré à l'encre, et un trait mène à
l'annotation, en haut du feuillet}\marginal{\uncertain{non}, pas tout à fait !}

\emph{Théorème} 7.1.\quad Soit $X \subset \mathbf{P}^r_k$ un schéma projectif
équidimensionnel de Cohen-Macaulay, de dimension $n > 0$, sur un corps $k$ de car.\
$p > 0$. Il existe alors un entier $d_0 > 0$ tel que, pour tout $d \geqslant d_0$,
si $Y$ est la section de $X$ par \struck{la} hypersurface générique de degré $d$,
définie sur l'extension $k' = k(t_0, \ldots, t_r)$ de $k$, on ait la propriété
qui suit. On désigne par $E_0(Y)$ le conoyau de $H^{n-1}(X_{k'},
\underline{\mathcal{O}}_{X_{k'}}) \to H^{n-1}(Y, \underline{\mathcal{O}}_Y)$, qui
est un vectoriel de dimension finie sur $k'$, muni de l'homomorphisme $p$-linéaire
de Frobenius $E_0(Y)^{(p)} \to E_0(Y)$ (provenant de la puissance $p$.ème sur les
faisceaux de coéfficients $\underline{\mathcal{O}}_{X_{k'}}$,
$\underline{\mathcal{O}}_Y$). Soit \add{$V =$} $E_0(Y)_{ss}$ la partie semi-simple
de $E_0(Y)$, i.e.\ le plus grand sous-espace \add{$V$} stable par frobenius et sur
lequel frobenius induise un morphisme surjectif $V^{(p)} \to V$. Considérons la
représentation correspondante de $G = \operatorname{Gal}(\bar k'/k')$ dans le
vectoriel \add{$V_0$} de dimension $s = \dim_{k'} V$ sur $\mathbf{F}_p$ formé des
$x \in V_{\bar k'}$, tels que $Fx = x$. Alors \struck{cette représentation est non
triviale, et de façon plus précise,} pour tout élément $a \in \mathbf{F}_p^{*}$,
il existe un $u \in \operatorname{Aut}(V_0)$ dans l'image \add{$G_0$} de $G$
\add{d'un sous-groupe $V_0' \subset V_0$ stable par $G_0$} tel que $\operatorname{Tr}
u|V_0' = a$. \struck{(Cela implique en particulier que $G_0$ n'est pas un
$p$-groupe, et plus précisément qu'il a au moins $p-1$ éléments distincts (et même
à traces distinctes) qui sont semi-simples.} \struck{Cela implique également que si
$\dim V \geqslant 2$).}\note{les ajouts « $V=$ », « $V$ », « $V_0$ », « $G_0$ »
sont à l'encre, dans l'interligne ; « d'un sous-groupe $V_0' \subset V_0$ stable
par $G_0$ » est une insertion à l'encre sous la ligne, dont la place dans la phrase
n'est pas marquée clairement, et le « $|V_0'$ » de la trace est ajouté à la main.
« et même » est tapé au-dessus d'un mot biffé. La parenthèse finale, avec sa suite
manuscrite (« Cela implique également que si $\dim V \geqslant 2$ »), est barrée
à l'encre et encadrée d'un crochet ; « $p-1$ » est tapé sur une autre lettre}

\marginal{\uncertain{pas prouvé sous cette forme juste}}\note{en marge gauche, à
l'encre, écrit de biais, en regard de l'énoncé du théorème, qu'un trait vertical
marque sur toute sa hauteur}

Signalons tout de suite les corollaires qui nous intéresseront surtout :

\emph{Corollaire} 7.2.\quad Supposons que $k$ soit fini. Soit $S_d$ le schéma
projectif défini par le vectoriel des formes de degré $d$ sur $\mathbf{P}^r_k$.
Alors pour tout ouvert non vide $U$ de $S_d$, et tout entier $a$, il existe un
\add{point} $t \in U(k')$ \add{($k'$ une extension finie de $k$)} tel
que\note{« point » remplace un mot tapé biffé à l'encre (« forme »,
semble-t-il) ; après « $U(k')$ », deux mots tapés sont barrés et surchargés d'une
lecture qui ne se laisse pas déchiffrer (\ill{}) ; « une extension finie de $k$ »
est écrit sous la ligne}

\marginal{\uncertain{pas prouvé ; également 7.1 sous sa forme juste}}\note{en
marge gauche, à l'encre, de biais, en regard du corollaire, qu'un trait vertical
marque aussi}

\page{18}

\note{p.~10 du tapuscrit}

\[ \operatorname{card}(Y_t(k')) \equiv a \mod p , \]
où $Y_t$ est la section de $X_{k'}$ par l'hypersurface de $\mathbf{P}^r_{k'}$ définie
par $t$.\note{dans la formule, « $(k')$ » remplace à l'encre une parenthèse tapée
et noircie ; de même « $k'$ » en indice de $X$ et de $\mathbf{P}^r$, à la place
d'un « par » et d'un indice tapés et barrés}

En effet, Deligne a vérifié, à l'aide de la formule de Verdier-Lefschetz, que pour un
schéma propre $Y$ sur un corps fini, la somme alternée des traces des frobenius
\add{géométriques} opérant sur $H^i(Y, \underline{\mathcal{O}}_Y)$ est égal au
nombre des points rationnels de $Y$ mod $p$.\note{« géométriques » est tapé dans
l'interligne et appelé à l'encre} D'autre part, quitte à restreindre $U$, nous
pouvons supposer que,
\[ f\colon \underline{Y} \longrightarrow S_d \]
étant le morphisme qu'on devine à fibres les $Y_s$ ($s \in S_d$), que \add{$f$ est
plat, et que} la formation des $R^i f_{*}(\underline{\mathcal{O}}_Y)$ commute au
changement de base au dessus de $U$, que les $R^i f_{*}(\underline{\mathcal{O}}_Y)$
sont localement libres sur $U$, enfin que $V$ est la fibre générique d'un sous-fibré
$\underline{V}$ de $R^{n-1} f_{*}(\underline{\mathcal{O}}_Y) /
\underline{\mathcal{O}}_{S_d} \otimes H^{n-1}(X, \underline{\mathcal{O}}_X)$,
stable par frobenius, et sur lequel frobenius soit semi-simple (i.e.\ un
épimorphisme $\underline{V}^{(p)} \to \underline{V}$). Alors la représentation du
groupe fondamental du point générique $s$ de $S_d$ envisagée dans 7.1 provient
d'une représentation du groupe fondamental $\pi$ de $U$, \struck{Quitte à
restreindre $U$, on peut supposer également que la partie semi-simple des fibres
génériques de tous les $R^i f(\underline{\mathcal{O}}_Y)$ (pas seulement pour
$i = n-1$) provient d'un sous-Module localement facteur direct, stable par
frobenius, sur lequel frobenius soit semi-simple.}\note{« $f$ est plat, et que »
est à l'encre au-dessus de la ligne ; le passage « Quitte à restreindre \ldots{}
semi-simple » est encadré et biffé à l'encre de grands traits obliques} D'autre part,
quitte à prendre $d_0$ assez grand dans 7.1, on peut supposer que $H^i(X,
\underline{\mathcal{O}}_X(-d)) = 0$ pour tout $i < n$ et tout $d \geqslant d_0$
(c'est ici qu'intervient l'hypothèse $X$ Cohen-Macaulay sur $X$), de sorte que pour
toute section $Y_s$ de $X$ de degré $d \geqslant d_0$ et de dimension $n-1$,
l'homomorphisme $H^i(X_s, \underline{\mathcal{O}}_{X_s}) \to H^i(Y_s,
\underline{\mathcal{O}}_{Y_s})$ sera un isomorphisme pour $i \leqslant n-2$,
\add{un monomorphisme pour $i = n-1$.}\note{l'ajout est à l'encre sous la ligne,
entouré ; un mot biffé à l'encre au milieu (« \ill{} ») ne se lit pas} On voit alors
que la formule de Lefschetz-Verdier donne, pour le point \struck{fermé} $t$ de
$U(k')$, que le nombre des points rationnels sur $k'$ de $Y_t$
satisfait\note{« le » corrige « un » à l'encre ; « fermé » est barré ; « $(k')$ »
et « $k'$ » remplacent, à l'encre, des lettres tapées ; « satisfait » est repassé à
la main}

\page{19}

\note{p.~11 du tapuscrit}

\[ \operatorname{card} Y_t(k') \equiv c_t + (-1)^{n-1} \operatorname{Tr} f_t , \]
où
\[ c_t = \sum_{0 \leqslant i \leqslant n-1} (-1)^i\, \text{Trace de }
\mathrm{Frob}_{k'} \text{ opérant sur } H^i(X, \underline{\mathcal{O}}_X) , \]
et où $f_t$ est l'élément du groupe \add{quotient} $G_0 =
\operatorname{Aut}_{\mathbf{F}_p}(V_0)$ de $\pi$, image de l'élément de frobenius
de $\pi$ associé au point \struck{fermé} $t \in U(k')$.\note{le signe $(-1)^i$ et
« quotient » sont à l'encre ; « image » est tapé après un mot biffé à la machine}
Quitte à \add{remplacer d'abord $k$ par} une extension finie, de sorte que
$\mathrm{Frob}_k$ opérant sur les $H^i(X, \underline{\mathcal{O}}_X)$ soit à
valeurs propres toutes égales à $0$ ou $1$, on peut supposer que $c_t$ est
indépendant de $t$, \add{soit $c$.} Il résulte alors de 7.1 qu'il existe un
élément $f$ de $G_0$ tel que $c + (-1)^{n-1} \operatorname{Tr} f \equiv a \mod p$,
et il suffit donc de choisir un $t \in U$ (considéré comme élément de $U(k(t))$) de
telle façon que $f_t = f$. Mais ceci est possible en vertu du théorème de
Čebotaref.\note{« remplacer d'abord $k$ par » et « soit $c$. » sont tapés dans
l'interligne et appelés à l'encre ; « passer à » est ainsi remplacé ; le caron de
Čebotaref est ajouté à la main}\marginal{pas prouvé}\note{en marge gauche, dans un
ovale au crayon, avec un trait vers la fin de la démonstration}

\emph{Corollaire} 7.3.\quad \add{(Si $n$ pair $\geqslant 4$, $X$ lisse).} On peut
trouver un ouvert $U'$ \add{(non vide de la grassmannienne des droites)} de $S_d$,
tel que pour toute droite \add{$D \in U'(k')$} \struck{de $(S_d)_{k'}$ qui
rencontre $U$ et qui définit} \add{$D$ est} un pinceau de Lefschetz sur
$X_{k'}$, si $t$ désigne le point générique de cette droite, le motif $E(Y_t)$ des
cycles évanescents sur $Y_t$ ne soit pas de niveau 1 (donc ne contienne pas de
sous-motif $\neq 0$ de niveau 1, pas même sur une extension de $k(t)$, d'après le
résultat de simplicité rappelé au n\textsuperscript{o}~6).\note{« (Si $n$ pair, $X$
lisse). » est tapé au-dessus de la ligne ; « $\geqslant 4$ », « (non vide de la
grassmannienne des droites) », « $D \in U'(k')$ » et « $D$ est » sont à l'encre ;
la fin de ligne « qui rencontre \ldots{} qui définit » est barrée à l'encre}
\marginal{($k'$ une ext.\ quelc.\ de $k$)}\note{en marge gauche, à l'encre, en
regard de « droite »}

\marginal{$U'$ est défini comme un \uncertain{ouvert} de droites $D$ rencontrant
$\overline{U}$ et telles que $\pi_1(D \cap U) \to \pi_1(\uncertain{U})$ soit
surjectif (\ill{} à $k'$). On établira un \uncertain{variante} de
Bertini}\note{note à l'encre écrite de bas en haut dans la marge gauche, encadrée,
et appelée par un trait devant « On peut »}

On prend $U$ comme dans la démonstration de 7.2, \add{et tel que les $Y_s$ ($s \in
U$) soient lisses.} On peut (grâce au fait qu'un motif de niveau $\leqslant 1$ se
spécialise en un motif de niveau $\leqslant 1$, comme on a signalé au
n\textsuperscript{o}~6 comme conséquence du critère de Ogg-Néron-Chafarévitch)
supposer que $k$ est fini, et il suffit de prouver que l'on peut trouver un point
fermé $s$ de $D$ tel que les valeurs propres de Frobenius opérant sur
$E^{n-1}(Y_s, \underline{\mathbf{Q}}_\ell)$ ($\ell$ un nombre premier fixé $\neq
p$) ne soient pas toutes divisibles par $p$. Or le nombre de points de $Y_t$ est
égal à la somme alternée des traces de frobenius opérant sur les\note{« et tel que
les $Y_s$ \ldots{} soient lisses. » est à l'encre au-dessus de la ligne ; un mot
tapé biffé suit « et il » ; le « $\neq \ell$ » du tapuscrit porte le caractère
$\ell$ pour $p$, et on lit $\neq p$}

\page{20}

\note{p.~12 du tapuscrit, le dernier du lot ; le bord droit du feuillet est coupé
sur la reproduction, et les fins de ligne y manquent}

$H^i(Y_{\bar t}, \underline{\mathbf{Q}}_\ell)$, qui peut s'écrire, grâce au
théorème de Lefschetz \uncertain{faible} \ldots\note{« faible » est à l'encre en
bout de ligne, en partie hors du cadre ; « aux théorèmes » est corrigé à la main
en « au théorème »}
\[ \sum_{0 \leqslant i \leqslant n-2} (-1)^i\, \text{Trace } \mathrm{Frob}_{H^i(X_{\bar
t}, \mathbf{Q}_\ell)} \bigl(1 + q^{\,2n \ldots}\bigr) + \ldots\, \text{Trace }
\mathrm{Frob}_{H^{n-1}(X_{\bar t}, \ldots)} \ldots \]
\[ \ldots + (-1)^{\ldots}\, \text{Trace } \mathrm{Frob}_{E^{n-1}(Y_s,
\underline{\mathbf{Q}}_\ell)} , \]
\marginal{$q = \operatorname{card} k(s)$}\note{la première ligne de la formule est
reprise à l'encre et n'est pas sûre : le facteur « $(1 + q^{2n\ldots})$ », écrit
au-dessus d'un « $(-1)^i$ » tapé et barré, a un exposant qu'on ne lit pas
(\ill{}) ; entre les deux traces, des signes $(-1)^{\ldots}$ successivement
biffés ; le signe $(-1)^i$ devant la somme et $(-1)^{\ldots}$ devant la dernière
trace sont à l'encre, sur un mot tapé biffé}

où tous les frobenius sont relatifs au corps de base $k(s)$. Ici encore, quitte à
remplacer $k$ par une extension finie, on peut supposer que les valeurs propres
\add{($0 \leqslant i \leqslant n-2$)} des frobenius (relatifs à $k$) opérant sur les
$H^i(X_{\bar k}, \underline{\mathbf{Q}}_\ell)$ \add{(en tant qu'entiers
algébriques)} sont égales à $0$ ou $1$ mod $p$, de sorte que mod $p$ la contribution
des termes de la première ligne soit indépendante du choix de $s$, \add{soit $c$.}
Donc si les valeurs propres de frobenius opérant sur les $E^{n-1}(Y_s,
\underline{\mathbf{Q}}_\ell)$ sont toutes divisibles par $p$, on en conclut que
$\operatorname{card} Y_s(k(t)) \equiv c \mod p$, et la même relation serait
également valable, pour la même raison, si on remplace $k(t)$ par une extension
finie $k'$ quelconque. Il suffit donc de voir qu'on peut trouver toujours un point
$s$ de $D \cap U$ \add{dans une extension finie $k'$ de $k$, tel que}
$\operatorname{card} Y_s(k')$ soit congru mod $p$ à un élément donné de
$\mathbf{F}_p$ (qu'on choisira alors $\neq c$). Or c'est ce que montre
essentiellement 7.2, à cela près qu'on n'a pas été assez précis dans l'énoncé de ce
corollaire puisqu'on peut imposer à $s$ de provenir de $D \cap U(k')$, où $D$ est
une droite donnée \add{$\in U'(k')$} \struck{rencontrant $U$}. Ce raffinement est
contenu dans la démonstration de 7.2, compte tenu du fait que par Bertini, le groupe
fondamental de $D \cap U_{k'}$ s'envoie \emph{sur} celui de $U$, donc s'envoie
\emph{sur} $G_0$.\note{« ($0 \leqslant i \leqslant n-2$) » et « (en tant
qu'entiers algébriques) » sont à l'encre ; « soit $c$. » et « dans une extension
finie $k'$ de $k$, tel que » sont tapés dans l'interligne et appelés à l'encre,
à la place de mots biffés à la machine (dont « fermé ») ; « $\in U'(k')$ » est à
l'encre et « rencontrant » barré à l'encre ; l'indice $k'$ de $U$ est à la main.
Le texte du tapuscrit s'arrête ici, au milieu du feuillet}

\end{document}
